\color{uniblue}\section[Описание и работа]{Описание и работа}
\color{black}

\color{uniblue}{\subsection{Назначение}}
\color{black}

\begin{enumerate}

\item
Устройство серии <<ЮНИТ-М300-ОЛ>> разработано для применения в схемах вторичной коммутации распределительных устройств 6-35 кВ с постоянным, переменным и выпрямленным оперативным током. Устройство предназначено для реализации функций релейной защиты отходящей линии напряжением 6-35 кВ, при этом может осуществлять функции измерения, регистрации, осциллографирования и сигнализации. При подключении устройства к информационной сети энергообъекта возможно осуществление функций дистанционного управления устройством и сбор данных с него посредством протоколов и интерфейсов, которые определяются конфигурацией устройства.
\item 
Устройство предлагается с типовой конфигурацией: в аппаратной части определяемой картой заказа (см. приложение \ref{app:card}), в функциональной части -- функционально-структурной схемой (см. приложение \ref{app:fsu}). Схемы подключения устройства по аналоговым входам приведены в приложении \ref{app:analogs}.
\item
Устройство серии <<ЮНИТ-М300-ОЛ>> разработано для установки в релейных отсеках ячеек КРУ, КРУН, камер КСО, а также для применения в составе шкафов для установки в релейных залах подстанций, построенных на базе архитектур I и II типа, а также на базе смешанной архитектуры.

\end{enumerate}

%==================== Технические характеристики устройства =============================

\color{uniblue}\subsection{Технические характеристики}
\color{black}

\begin{enumerate}

\item
Основные номинальные параметры устройства указаны в таблице \ref{desc:chars}. 

\setlength{\extrarowheight}{0.06cm} % высота ячеек в таблице
\small
\begin{longtable}
    {|>{\centering\arraybackslash}m{12.9cm}|>{\centering\arraybackslash}m{4cm}|}
    \caption{Основные номинальные параметры устройства\hfill\vspace{-0.5\baselineskip}}\label{desc:chars}\\ 
    \hhline{|-|-|}
    \rowcolor{gray!30}
    Параметр & Значение \\
    \hhline{|-|-|}
    \endfirsthead

    %\caption{\emph{Продолжение\hfill\vspace{-0.5\baselineskip}}} \\ % Отступ обозначения таблицы от таблицы и выравнивание надписи слева
    \caption*{\hspace{3pt}\emph{Продолжение таблицы \ref{desc:chars}\hfill\vspace{-0.5\baselineskip}}} \\ % сделано по ГОСТ 2.105 п.6.8.7
    \hhline{|-|-|}
    \rowcolor{gray!30}
    Параметр & Значение \\ 
    \hhline{|-|-|}
    \endhead
    \hhline{|-|-|}

    \endfoot
    \endlastfoot

    \raggedright Номинальный переменный ток, А  & \raggedright\arraybackslash 0,2; 1; 5 \\  \hhline{|-|-|}
    \raggedright Номинальный переменное фазное напряжение, В  & \raggedright\arraybackslash 100/$\sqrt{3}$ \\  \hhline{|-|-|}
    \raggedright Номинальная частота, Гц  & \raggedright\arraybackslash 50 \\  \hhline{|-|-|}
    \raggedright Номинальное напряжение оперативного постоянного тока, В  & \raggedright\arraybackslash 110; 220 \\  \hhline{|-|-|}
    \raggedright Номинальное напряжение оперативного переменного тока, В  & \raggedright\arraybackslash 110; 220 \\  \hhline{|-|-|}

\end{longtable}
\normalsize
\setlength{\extrarowheight}{0.0cm} % возврат параметра высота ячеек в таблице

\end{enumerate}
%================================================================ Конструктивное исполнение ==================================================

\color{uniblue}\subsection{Конструктивное исполнение}
\color{black}

\begin{enumerate}

\item
Корпус устройства в типовой конфигурации предусматривает отдельные исполнения по ширине в зависимости от требуемого количества блоков дискретного управления и дискретных входов, габаритные размеры устройства приведены в приложении \ref{app:sizes}. Внешний вид устройства для различных типовых исполнений показан на рисунках \ref{device:img1}, \ref{device:img2}. Подробное описание конструкции и функций используемых блоков (а также доступных к заказу) приведено в документации ЮТКБ.656122.609 РЭ <<Устройство защиты и автоматики серии <<ЮНИТ-М300>>. Руководство по эксплуатации общее>>.

\item
По желанию Заказчика возможно выполнение устройства в других конструктивных исполнениях в соответствии с общей картой заказа на серию <<ЮНИТ-М300>>, которая приводится в документации ЮТКБ.656122.609 РЭ <<Устройство защиты и автоматики серии <<ЮНИТ-М300>>. Руководство по эксплуатации общее>>. 

\item
Для типоисполнения <<ЮНИТ-М314>> предусматривается установка блоков в любой из слотов с <<М3>> по <<М7>>:
\begin{itemize}
    \item выходных реле типа <<К001>> или <<К002>>;
    \item дискретных входов типа <<В001>> или <<В021>>.
\end{itemize}
\item
Для типоисполнения «ЮНИТ-М319» предусматривается установка блоков в любой из слотов с <<М3>> по <<М11>>:
\begin{itemize}
    \item выходных реле типа <<К001>>, <<К002>>;
    \item дискретных входов типа <<В001>> или <<В021>>.
\end{itemize}

\vspace{3mm}
\begin{figure}[H]
\centering
\includegraphics[width=0.7\textwidth,height=0.7\textheight,keepaspectratio]{314.jpg}
\captionof{figure}{Вид устройства <<ЮНИТ-М314-ОЛ>> с размещением блоков по слотам}\label{device:img1}
\end{figure}

\vspace{3mm}
\begin{figure}[H]
\centering
\includegraphics[width=0.9\textwidth,height=0.9\textheight,keepaspectratio]{319.jpg}
\captionof{figure}{Вид устройства <<ЮНИТ-М319-ОЛ>> с размещением блоков по слотам}\label{device:img2}
\end{figure}

\item
В устройстве устанавливается измерительный блок в зависимости от карты заказа, например блок <<М044>> (в котором предусмотрено четыре канала для измерения переменного тока и четыре канала для измерения переменного напряжения), в устройстве <<ЮНИТ-М314>> занимает место в слотах <<М08>>, <<М09>>, а в устройстве <<ЮНИТ-М319>> занимает место в слотах <<М12>>, <<М13>>. Типовое подключение входных аналоговых сигналов к измерительному блоку в составе устройства показано в приложении \ref{app:analogs}.

\item
Подробное описание параметров настроек и технических характеристик блоков дискретных входов, блоков дискретного управления и измерительных блоков приводится в документации ЮТКБ.656122.609 РЭ <<Устройство защиты и автоматики серии <<ЮНИТ-М300>>. Руководство по эксплуатации общее>>.

\end{enumerate}

%================================================================ Функциональный состав  ==================================================
\color{uniblue}\subsection{Функциональный состав}
\color{black}

\begin{enumerate}

\item
Перечень функций, реализованных в составе устройства, приведен в таблице \ref{desc:tbl1}. Описание функций релейной защиты и автоматики в составе устройства, а также их функциональные схемы и уставки приводятся в разделе \ref{sec:funcs}. Подробное описание вспомогательных (сервисных) функций приведено в документации ЮТКБ.656122.609 РЭ <<Устройство защиты и автоматики серии <<ЮНИТ-М300>>. Руководство по эксплуатации общее>>. Общая функционально-структурная схема взаимосвязей функций, дискретных входов и выходов устройства приведена в приложении \ref{app:fsu}. Принятые обозначения элементов схем приведены в приложении \ref{app:func_desc}.

\setlength{\extrarowheight}{0.05cm} % высота ячеек в таблице
\small
\begin{longtable}
{|>{\centering\arraybackslash}m{3.2cm}|>{\centering\arraybackslash}m{13.7cm}|}
\caption{Функции в составе устройства <<ЮНИТ-М300-ОЛ>>\hfill\vspace{-0.5\baselineskip}}\label{desc:tbl1}\\ 
\hhline{|-|-|}
\rowcolor{gray!30}
Группа функций & Примечание \\ 
\hhline{|-|-|}
\endfirsthead

\caption*{\hspace{3pt}\emph{Продолжение таблицы \ref{desc:tbl1}\hfill\vspace{-0.5\baselineskip}}} \\ % сделано по ГОСТ 2.105 п.6.8.7
\hhline{|-|-|}
\rowcolor{gray!30}
Группа функций & Примечание \\ 
\hhline{|-|-|}
\endhead
\hhline{|-|-|}

\endfoot
\endlastfoot

\centering Основные функции & \raggedright\arraybackslash Функции защиты отходящего присоединения 6-35 кВ для различных режимов заземления нейтрали (ДЗ, ТО, МТЗ, ТЗНП, ЗОП, ЗП, ГСОЗЗ, НЗОЗЗ и пр.), функции автоматики отходящего присоединения 6-35 кВ (АПВ, КП, управление коммутационными аппаратами), УРОВ  \\  \hline

\centering Сервисные функции & \raggedright\arraybackslash Сигнализация (обеспечение ввода служебных дискретных сигналов о состоянии оборудования шкафа РЗА, формирование отчетов по изменению состояния дискретных внешних и внутренних сигналов и выдачу отчетов в АСУ~ТП, формирование и выдачу сигналов аварийной и предупредительной сигнализации в схему резервной сигнализации) \\ \cline{2-2}
& \raggedright\arraybackslash  Осциллографирование и журналирование (определение места повреждения, запись доаварийных, аварийных и послеаварийных режимов с сохранением осциллограмм в энергонезависимой памяти устройства, запись о событиях по изменению конфигурации устройства, изменения состояния внутренних сигналов алгоритмов, обновление встраиваемого программного обеспечения и пр. с метками времени и описанием события в энергонезависимой памяти устройства) \\ \cline{2-2}
& \raggedright\arraybackslash Самодиагностика (непрерывный оперативный контроль работоспособности в течение всего времени работ) \\ \cline{2-2}
& \raggedright\arraybackslash Конфигурирование (задание/переключение уставок, ввод в работу / вывод из работы, изменение параметров самоописания устройства, перезагрузка устройства, ввод/вывод тестового режима, изменение режима работы информационных сервисов, возможность переопределения связей входных и выходных логических сигналов алгоритмов устройства с дискретными входами, выходными реле, светодиодными индикаторами и функциональными клавишами управления) \\ \cline{2-2}
& \raggedright\arraybackslash Тестовый режим (обработка принятых сигналов и отправка отчетов с установленным флагом <<Тест>>) \\ \cline{2-2}
& \raggedright\arraybackslash Самоописание (предоставление данных об устройстве посредством устройства <<ЮНИТ-ИЧМ>> и сервисную программное обеспечение <<ЮНИТ-СЕРВИС>>) \\ \cline{2-2}
& \raggedright\arraybackslash Синхронизация времени (календарь и часы реального времени, синхронизируемые от персонального компьютера или от АСУ~ТП) \\ \hhline{|-|-|}

\end{longtable}
\setlength{\extrarowheight}{0.0cm} % возврат параметра высота ячеек в таблице
\normalsize

\item
Устройство поддерживает два режима оперативного управления: местный и дистанционный. Выбор режима управления устройством производится посредством функциональной клавиши <<М/Д>> или по состоянию выбранного дискретного входа. В местном режиме игнорируются все команды управления с верхнего уровня АСУ ТП. Более подробно режимы устройства описаны в документации ЮТКБ.656122.609 РЭ <<Устройство защиты и автоматики серии <<ЮНИТ-М300>>. Руководство по эксплуатации общее>>.
\item
Для управления режимами работы основных функций (оперативный ввод/вывод, перевод на сигнал и пр.), а также для генерации различных сигналов сброса функций (сброс триггеров фиксации неисправности цепей газовых и технологических защит, сброс счетчиков и пр.), применяются виртуальные ключи и кнопки. Подробное описание и сценарии работы виртуальных ключей и кнопок приведено в документации ЮТКБ.656122.609 РЭ <<Устройство защиты и автоматики серии <<ЮНИТ-М300>>. Руководство по эксплуатации общее>>. 

%%%%%%%%%%%%%%%%%%%%%%%%%%%%%% Группы уставок %%%%%%%%%%%%%%%%%%%%%%%%%%%%%%
\item Группы уставок
\begin{enumerate}

\item Устройство имеет четыре группы уставок, предусмотрено оперативное переключение между несколькими сохраненными группами уставок. Подробное описание работы с группами уставок приведено в документации ЮТКБ.656122.609 РЭ <<Устройство защиты и автоматики серии <<ЮНИТ-М300>>. Руководство по эксплуатации общее>>.

\item
Изменение группы уставок возможно:
\begin{itemize}
    \item по месту, через устройство <<ЮНИТ-ИЧМ>>, п.\ref{sec1:hmi};
    \item дистанционно, через команды от АСУ ТП;
    \item дистанционно, средствами сервисного ПО <<ЮНИТ-СЕРВИС>>;
    \item по месту, через соответствующий дискретный вход сигналом импульсного типа.
\end{itemize}
\item
Переключение группы уставок происходит не мгновенно. После появления сигнала/команды на изменение активной группы уставок и до момента активации новой группы проходит от 7 до 15 с. В это
время происходит проверка целостности файлов выбранной группы уставок. Всё это время устройство продолжает функционировать на той группе уставок, которая была активна до появления сигнала/команды на изменение активной группы уставок.

\end{enumerate}
%%%%%%%%%%%%%%%%%%%%%%%%%%%%%%%%%%%%%%%%%%%%%%%%%%%%%%%%%%%%

%%%%%%%%%%%%%%%%%%%%%%%%%%%%%% Аварийный осциллограф %%%%%%%%%%%%%%%%%%%%%%%%%%%%%%
\item Аварийный осциллограф
\begin{enumerate}
\item
При возникновении нарушений нормальных режимов работы данный процесс может записываться в память с помощью аварийного осциллографа. Осциллографирование выполняется с частотой дискретизации не менее 1200 Гц. Осциллограммы содержат измеренные значения аналоговых сигналов, дискретные сигналы состояния входов, выходных реле, а также логических сигналов функций РЗА и сервисные сигналы устройства. Подробное описание параметров встроенного регистратора аварийных событий приведено в документации ЮТКБ.656122.609 РЭ <<Устройство защиты и автоматики серии <<ЮНИТ-М300>>. Руководство по эксплуатации общее>>.
\item
Пуск осциллографа производится по факту изменения состояния дискретных сигналов, заданных в качестве критериев пуска. В зависимости от длительности записи предаварийного и послеаварийного режима, в памяти может храниться разное количество осциллограмм. Данные, собранные аварийным осциллографом, сохраняются в энергонезависимой памяти устройства и циклически перезаписываются. Удаление записей аварийного осциллографа пользователем невозможно.

\item
Перечень дискретных сигналов, доступных для регистрации аварийным осциллографом, приведен в приложении \ref{app:signals}. Перечень аналоговых сигналов, доступных для регистрации аварийным осциллографом, приведен в таблице \ref{osz:tbl1}.


\setlength{\extrarowheight}{0.05cm}
\small
\begin{longtable}{|>{\centering\arraybackslash}m{6.8cm}|>{\centering\arraybackslash}m{2cm}|>{\centering\arraybackslash}m{5.3cm}|>{\centering\arraybackslash}m{2cm}|}
\caption{Перечень аналоговых сигналов, доступных для регистрации аварийным осциллографом\hfill\vspace{-0.5\baselineskip}}\label{osz:tbl1}\\ 
\hline
\rowcolor{gray!30}
Описание аналогового сигнала & Канал & Наименование  & Пуск РАС \\ 
\hline
\endfirsthead
\caption*{\hspace{3pt}\emph{Продолжение таблицы \ref{osz:tbl1}\hfill\vspace{-0.5\baselineskip}}} \\ % сделано по ГОСТ 2.105 п.6.8.7
\hline
\rowcolor{gray!30}
Описание аналогового сигнала & Канал & Наименование  & Пуск РАС \\ 
%\hline
\endhead
%\hline
\endfoot
%\hline
\endlastfoot

\raggedright Фазные токи & \centering 1 & \centering Ia  & \centering\arraybackslash + \\
\hline
\raggedright >>             & \centering 2 & \centering Ib  & \centering\arraybackslash + \\
\hline
\raggedright >>             & \centering 3 & \centering Ic  & \centering\arraybackslash + \\
\hline

\raggedright Фазные напряжения & \centering 4 & \centering Ua  & \centering\arraybackslash + \\
\hline
\raggedright >>              & \centering 5 & \centering Ub  & \centering\arraybackslash + \\
\hline
\raggedright >>              & \centering 6 & \centering Uc  & \centering\arraybackslash + \\
\hline

\raggedright Напряжение <<разомкнутого треугольника>>  & \centering 7 & \centering Uhk  & \centering\arraybackslash + \\
\hline

\raggedright Частота & \centering 8 & \centering F  & \centering\arraybackslash + \\
\hline

\raggedright Вычисленный ток ОП & \centering 9 & \centering I2  & \centering\arraybackslash + \\
\hline
\raggedright Вычисленный ток НП & \centering 10 & \centering 3I0  & \centering\arraybackslash + \\
\hline

\raggedright Вычисленное напряжение ОП & \centering 11 & \centering U2  & \centering\arraybackslash + \\
\hline
\raggedright Вычисленный напряжение НП & \centering 12 & \centering 3U0  & \centering\arraybackslash + \\
\hline

\end{longtable}
\normalsize
\setlength{\extrarowheight}{0.0cm}

\end{enumerate}
%%%%%%%%%%%%%%%%%%%%%%%%%%%%%%%%%%%%%%%%%%%%%%%%%%%%%%%%%%%%

%%%%%%%%%%%%%%%%%%% Журнал событий %%%%%%%%%%%%%%%%%%%%%%%%%
\item Журнал событий
\begin{enumerate}

\item
В устройстве реализована функция записи сигналов дискретных входов, работы защит, автоматики и сервисных сигналов в журнал событий. События безопасности записываются в отдельный журнал безопасности. Подробное описание параметров журналов событий приведено в документации ЮТКБ.656122.609 РЭ <<Устройство защиты и автоматики серии <<ЮНИТ-М300>>. Руководство по эксплуатации общее>>.
\item
Максимальное число записей журнала -- 20000. Каждая запись в журнал событий содержит в себе следующую информацию:
\begin{itemize}
\item дату и время возникновения события;
\item идентификатор события;
\item состояние дискретного или значение аналогового сигнала в момент события;
\item качество зарегистрированного параметра;
\item качество времени.
\end{itemize}
\item
Данные о событиях сохраняются в энергонезависимой памяти. При заполнении журнала событий перезапись осуществляется от более старых событий к новым, то есть циклически. Удаление информации из журнала событий пользователем невозможно. Перечень сигналов, доступных для регистрации в журнале событий, приведен в приложении \ref{app:signals}.

\end{enumerate}
%%%%%%%%%%%%%%%%%%%%%%%%%%%%%%%%%%%%%%%%%%%%%%%%%%%%%%%%%%%%

%%%%%%%%%%%%%%%%%%%%%%%%%%%%%%Режим тестирования %%%%%%%%%%%%%%%%%%%%%%%%%%%%%%
\item Режим тестирования
\begin{enumerate}

\item В устройстве реализована поддержка режима тестирования <<Тест>> в объеме требований стандартов МЭК 61850. Режим тестирования предусмотрен для возможности оперативной проверки элементов/функций, таких как: ФСУ, измерительные органы, дискретные входы и выходные реле устройства.

\item  Особенности перевода устройства в режим тестирования и функциональности устройства в этом режиме описаны в документации ЮТКБ.656122.609 РЭ <<Устройство защиты и автоматики серии <<ЮНИТ-М300>>. Руководство по эксплуатации общее>>.

\end{enumerate}
%%%%%%%%%%%%%%%%%%%%%%%%%%%%%%%%%%%%%%%%%%%%%%%%%%%%%%%%%%%%

%%%%%%%%%%%%%%%%%%%%%%%%%%%%%% Синхронизация времени %%%%%%%%%%%%%%%%%%%%%%%%%%%%%%
\item Синхронизация времени
\begin{enumerate}

\item
Устройство предусматривает возможность синхронизации внутренних часов реального времени несколькими способами:
\begin{itemize}
\item по протоколу SNTP;
\item по протоколу PTP;
\item через дискретный вход устройства с использованием сигнала синхронизации 1PPS;
\item {средствами стандартизованного протокола передачи данных.}
\end{itemize}
Подробное описание способов синхронизации времени устройства приведено в документации ЮТКБ.656122.609 РЭ <<Устройство защиты и автоматики серии <<ЮНИТ-М300>>. Руководство по эксплуатации общее>>.

\end{enumerate}
%%%%%%%%%%%%%%%%%%%%%%%%%%%%%%%%%%%%%%%%%%%%%%%%%%%%%%%%%%%%

%%%%%%%%%%%%%%%%%%%%%%%%%%%%%% Система самодиагностики %%%%%%%%%%%%%%%%%%%%%%%%%%%%%%
\item Система самодиагностики
\begin{enumerate}

\item
Система самодиагностики устройства выполняет тесты в полном объеме при включении устройства, а также непрерывно в фоновом режиме.
Системой самодиагностики устройства в автоматическом режиме контролируются следующие параметры:
\begin{itemize}
\item состояние аппаратной части устройства, в том числе: АЦП измерительных блоков, блока питания, ОЗУ, ПЗУ, процессора, блоков дискретных входов, блоков дискретного управления, вспомогательных блоков;
\item температурный режим устройства;
\item состояние синхронизации времени;
\item сохранность исполняемого программного кода (целостность ПО);
\item состояние измерительных цепей;
\item исправность используемых каналов связи.
\end{itemize}

\item
Система самодиагностики фиксирует критические неисправности (<<Неисправность>>) и некритические ошибки (<<Ошибка>>). По факту выхода из режимов критической неисправности и предупредительной сигнализации производится запись в журнале событий. Подробное описание системы самодиагностики приводится в документации ЮТКБ.656122.609 РЭ <<Устройство защиты и автоматики серии <<ЮНИТ-М300>>. Руководство по эксплуатации общее>>.

\item
В случае появления критической неисправности активируется внутренний сигнал <<Критическая неисправность>>, производится запись в журнал событий, замыкаются контакты реле IRF, прекращается работа алгоритмов основных функций в составе устройства, это состояние устройства сигнализируется свечением светодиода <<Неисправность>> и отсутствием свечения светодиода <<Готовность>>. Для выхода из режима критической неисправности требуется перезапуск устройства.

\item
В случае появления некритической ошибки активируется внутренний сигнал <<Некритическая ошибка>>, производится запись в журнал событий, алгоритмы основных функций в составе устройства продолжают выполняться, это состояние устройства сигнализируется свечением светодиода <<Неисправность>> и отсутствием свечения светодиода <<Готовность>>. Устройство может выйти из режима предупредительной сигнализации без активных действий со стороны пользователя по факту пропадания критериев режима.

\end{enumerate}
%%%%%%%%%%%%%%%%%%%%%%%%%%%%%%%%%%%%%%%%%%%%%%%%%%%%%%%%%%%%

%%%%%%%%%%%%%%%%%%%%%%%%%%%%%% Аутентификация и авторизация пользователей %%%%%%%%%%%%%%%%%%%%%%%%%%%%%%
\item Аутентификация и авторизация пользователей
\begin{enumerate}

\item Устройство имеет встроенную реализацию правил разграничения прав доступа, для этого в устройстве существует ролевая модель разграничения прав доступа.
В устройстве предусмотрены следующие роли:
\begin{itemize}
\item Гость (пароль не требуется, доступен просмотр данных устройства);
\item Дежурный;
\item Инженер;
\item Офицер (сотрудник) безопасности;
\item Администратор (пароль <<АААА>> (кириллица), учетная запись предназначена для первоначального конфигурирования штатного профиля пользователей).
\end{itemize}

\item Все пользователи делятся на непривилегированных (роль Гость) и привилегированных (все остальные роли).

\item
Предустановленные пользователи с именами Гость и Администратор не подлежат удалению и принудительной блокировке, их роли не подлежат изменению. Гость используется как роль по умолчанию для доступа к просмотру данных устройства в режиме чтение без запроса аутентификационных данных пользователя. Администратор необходим для добавления/удаления других пользователей, установки им соответствующих ролей и установке/изменении паролей пользователей.
\item
После перезагрузки устройства подключенная сессия привилегированного пользователя сбрасывается. Устройство запускается с сеансом непривилегированного пользователя (роль Гость).
Подробное описание ролей и доступных полномочий и ограничений для каждой из ролей приведено в документации ЮТКБ.656122.609 РЭ <<Устройство защиты и автоматики серии <<ЮНИТ-М300>>. Руководство по эксплуатации общее>>.

\item
Для возможности работы с устройством пользователь должен пройти аутентификацию на устройстве. Аутентификация проводится по имени и паролю, указанными пользователем. 
\item
Функции безопасности устройства предъявляют следующие требования к паролям пользователей:
\begin{itemize}
    \item пароль может содержать буквы и/или цифры;
    \item длина пароля не может быть меньше семи символов. 
\end{itemize}
\item
Хранение паролей осуществляется в памяти устройства в нечитаемом виде, не позволяющем восстановить пароль.
\item
Функции безопасности не дают возможность повторного применения ранее использованных паролей. По умолчанию глубина хранения ранее заданных пользователем паролей в устройстве составляет четыре пароля.
\item
В процессе ввода пароля пользователю отображается только вводимый символ, ранее введенные символы не отображаются в открытом виде.
\item
Функции безопасности устройства ограничивают количество неуспешных попыток аутентификации за установленный период времени. Количество попыток можно задать в диапазоне от 1 до 9 (по умолчанию – 3)
\item
В случае превышения количества неуспешных попыток аутентификации:
\begin{itemize}
    \item выполняется блокировка доступа пользователя к устройству. Длительность блокировки определяется в настройках устройства пользователем с ролью Администратор в диапазоне от 1 минуты до 9999 минут (по умолчанию – 30 минут);
    \item формируется соответствующее событие журнала событий безопасности;
    \item формируется сигнализация для АСУ ТП о превышении заданного количества неуспешных попыток доступа. 
\end{itemize}
\item
В случае, когда количество неуспешных попыток аутентификации пользователя больше нуля, но меньше максимально допустимого, оно будет сброшено для данного пользователя по истечении времени сброса неуспешных попыток (5 минут).
\item
Блокировка возможности аутентификации пользователя снимается в случае:
\begin{itemize}
    \item истечения времени блокировки входа;
    \item принудительной разблокировки пользователем с ролью Администратор. 
\end{itemize}

\item
Пользователь может быть заблокирован принудительно пользователем с ролью Администратор. Разблокирование данного пользователя может осуществить так же только пользователь с ролью Администратор.
\item
Все настраиваемые параметры безопасности (время неактивности пользователя, количество безуспешных попыток входа, время блокировки входа) разрешены для редактирования пользователям с ролью офицера безопасности. При сбросе устройства к заводским настройкам, параметры безопасности сохраняются.

%\noindent
%\colorbox{yellow!40}{
    %\begin{minipage}{\dimexpr\linewidth-2\fboxsep}
    %{\faExclamationTriangle} При утере пароля учетной записи пользователя с ролью Администратор, восстановление административных прав возможно только путем перепрограммирования устройства специалистом предприятия-изготовителя, что повлечет за собой сброс настроечных параметров на заводские.
    %\end{minipage}
%}

% ПРЕДУПРЕЖДЕНИЕ - Несоблюдение указанных требований может нанести ущерб оборудованию или нарушению работы энергосистемы
%\noindent
%\fboxrule=1.2pt % задаём толщину рамки (жирная линия)
%\noindent
%\fbox{%
    %\begin{minipage}{\dimexpr\linewidth-2\fboxrule-2\fboxsep}
    %{\hspace*{0.5em}\textcolor{unired}{\Large\faExclamationTriangle}}\hspace{0.5em}
    %При утере пароля учетной записи пользователя с ролью Администратор, восстановление административных прав возможно только путем перепрограммирования устройства специалистом предприятия-изготовителя, что повлечет за собой сброс настроечных параметров на заводские.
    %\end{minipage}%
%}

% ОПАСНОСТЬ - Несоблюдение указанных требований может повлечь угрозу здоровью и жизни
%\noindent
%\fboxrule=1.2pt % задаём толщину рамки (жирная линия)
%\noindent
%\fbox{%
    %\begin{minipage}{\dimexpr\linewidth-2\fboxrule-2\fboxsep}
    %{\hspace*{0.5em}\textcolor{red}{\Large\faBolt}}\hspace{0.5em}
    %При утере пароля учетной записи пользователя с ролью Администратор, восстановление административных прав возможно только путем перепрограммирования устройства специалистом предприятия-изготовителя, что повлечет за собой сброс настроечных параметров на заводские.
    %\end{minipage}%
%}

% ИНФОРМАЦИЯ - 
\par
\noindent
\begin{tabularx}{\linewidth}{p{2em} X}
  \begin{minipage}[c][3\baselineskip][b]{2em}
    \centering\textcolor{uniblue}{\Huge\faCircleInfo}
  \end{minipage} &
  \hyphenpenalty=0 \textbf{При утере пароля учетной записи пользователя с ролью Администратор, восстановление административных прав возможно только путем перепрограммирования устройства специалистом предприятия-изготовителя, что повлечет за собой сброс настроечных параметров на заводские.}
\end{tabularx}
\vskip 2mm

\end{enumerate}
%%%%%%%%%%%%%%%%%%%%%%%%%%%%%%%%%%%%%%%%%%%%%%%%%%%%%%%%%%%%

%%%%%%%%%%%%%%% Блок интерфейсный (ЮНИТ-ИЧМ) %%%%%%%%%%%%%%%
\item\label{sec1:hmi} Блок интерфейсный (ЮНИТ-ИЧМ)

\begin{enumerate}

\item
Для местного управления и контроля за состоянием устройства предусматривается подключение блока интерфейсного <<ЮНИТ-ИЧМ>> производства ООО <<Юнител Инжиниринг>>. Блок подключается к <<ЮНИТ-М300>> посредством патч-корда с разъемами RJ45. В устройстве <<ЮНИТ-М300>> для подключения к блоку может использоваться отдельный порт X5, расположенный на блоке центрального процессора или любой из задействованных портов связи Ethernet (X1, Х2, Х3, X4).
\item
Более подробная информация по работе с устройством <<ЮНИТ-ИЧМ>> приведена в документации RU.37182817.00001.02.34.01 <<Программное обеспечение мониторинга и параметрирования <<ЮНИТ-СЕРВИС>>. Руководство пользователя>>.

\item
Подробное описание блока приведено в документе ЮТКБ.656132.701 РЭ – <<Блок интерфейсный <<ЮНИТ-ИЧМ>>. Руководство по эксплуатации>>. Блок интерфейсный <<ЮНИТ-ИЧМ>> является опциональным для заказа.

\end{enumerate}
%%%%%%%%%%%%%%%%%%%%%%%%%%%%%%%%%%%%%%%%%%%%%%%%%%%%%%%%%%%%

\end{enumerate}

%============================================================ Организация обмена информацией ==================================================
\color{uniblue}\subsection{Организация обмена информацией}\color{black}

\begin{enumerate}

\item\sloppy Описание и настройки доступных портов и интерфейсов связи приводится в документации ЮТКБ.656122.609 РЭ <<Устройство защиты и автоматики серии <<ЮНИТ-М300>>. Руководство по эксплуатации общее>>.

\item
В устройствах с поддержкой протокола передачи данных МЭК 61850 в качестве входных сигналов ФСУ могут быть использованы данные из GOOSE-сообщений, на которые оформлена подписка. Количество сигналов GOOSE, получаемых от интеллектуальных устройств, в данном типоисполнении устройства равно 28 шт.
\item
Сигналы из GOOSE-сообщений проходят предварительную обработку перед тем как попасть в <<Блок объединительный>> и далее в ФСУ (см. приложение \ref{app:fsu}).
В предварительную обработку входят следующие операции:
\begin{itemize}
\item проверка качества принимаемого сигнала;
\item алгоритм обработки резервирования GOOSE-сообщений;
\item алгоритм обработки симуляции GOOSE-сообщений;
\item алгоритм обработки метки <<Test>> GOOSE-сообщений;
\item подстановка значения по умолчанию при <<плохом>> качестве данных, принимаемых из GOOSE-сообщений.
\end{itemize}
Подробнее о подписке на GOOSE-сообщения указано в RU.37182817.00001.02.34.01 – <<Программное обеспечение мониторинга и параметрирования <<ЮНИТ-СЕРВИС>>. Руководство пользователя>>.

\item 
В устройствах с поддержкой протокола передачи данных МЭК 61850 для связи со смежными устройствами может быть использован протокол GOOSE по МЭК 61850-8.1. Количество сигналов GOOSE-сообщений, передаваемых в смежные интеллектуальные устройства, зависит от типоисполнения ЮНИТ-М300, первичного оборудования и в общем случае составляет не менее 100 шт. Подробнее о публикации GOOSE-сообщений указано в RU.37182817.00001.02.34.01 – <<Программное обеспечение мониторинга и параметрирования <<ЮНИТ-СЕРВИС>>. Руководство пользователя>>.

\end{enumerate}



%============================================================ Средства измерения, инструмент и принадлежности ==================================================
\color{uniblue}\subsection{Средства измерения, инструмент и принадлежности}\color{black}
Перечень оборудования и средств измерения, рекомендуемый для проведения эксплуатационных проверок устройства, приведен в документации ЮТКБ.656122.609 РЭ <<Устройство защиты и автоматики серии <<ЮНИТ-М300>>. Руководство по эксплуатации общее>>.

%================================================================= Маркировка и пломбирование ==================================================
\color{uniblue}\subsection{Маркировка и пломбирование}\color{black}
Сведения по маркировке и пломбировании устройства приведены в документации ЮТКБ.656122.609 РЭ <<Устройство защиты и автоматики серии <<ЮНИТ-М300>>. Руководство по эксплуатации общее>>.

%======================================================================== Упаковка ==================================================
\color{uniblue}\subsection{Упаковка}\color{black}
Сведения по упаковке устройства приведены в документации ЮТКБ.656122.609 РЭ <<Устройство защиты и автоматики серии <<ЮНИТ-М300>>. Руководство по эксплуатации общее>>.
